\section{The Momentum Transformer: An Intelligent Architecture}

\begin{frame}{The Design Philosophy: Story over Statistics}
    Why a Transformer? Momentum trading is often a battle against \textbf{Non-Stationarity}.
    \begin{itemize}
        \item Linear models assume tomorrow looks like today.
        \item LSTMs assume the recent past is all that matters.
        \item \textbf{The Transformer} treats history as a searchable database of regimes.
    \end{itemize}
    \textit{The Storyline:} We move from filtering noise (VSN), to capturing local context (LSTM), to cross-referencing multi-year regimes (Self-Attention).
\end{frame}

\begin{frame}{Variable Selection Networks (VSN): Separating Signal from Noise}
    The VSN acts as a "Gated Filter" for each asset:
    \begin{equation}
        \tilde{x}_t = \sum_{i=1}^m v_t^{(i)} \cdot \tilde{x}_t^{(i)}
    \end{equation}
    \begin{itemize}
        \item Automatically identifies which MACD scales (10-day vs 252-day) are dominant in the current regime.
        \item In high-noise environments (e.g., sideways equity markets), the VSN dampens uninformative features to prevent overfitting.
    \end{itemize}
\end{frame}

\begin{frame}{The Architecture: Hybridizing Recurrence and Attention}
    \begin{itemize}
        \item \textbf{LSTM Layer:} Processes local, path-dependent patterns. It handles the "short-term memory" and local volatility clustering.
        \item \textbf{Multi-Head Attention:} Allows the model to jump back in time. If today's sell-off looks like 2008, the model "attends" to that regime's recovery pattern.
        \item \textbf{Residual "Handcuffs":} We anchor the output to a proven TSMOM prior. The model only learns the \textit{correction} factor, not the trade itself.
    \end{itemize}
\end{frame}

\begin{frame}{High-Fidelity Engineering: Scaling to 50 Assets}
    Replicating Wood et al. (2022) required a significant jump in complexity:
    \begin{itemize}
        \item \textbf{Asset Count:} 24 $\to$ 50 (Full breadth of global macro futures).
        \item \textbf{Horizon:} 64 days $\to$ 252 days (Capturing annual cycles).
        \item \textbf{The Memory Bottleneck:} Long sequences in financial data create exponential VRAM requirements.
    \end{itemize}
\end{frame}

\begin{frame}{Technical Post-Mortem: Hitting the CUDA Wall}
    Increasing sequence length to 252 days triggered immediate **CUDA Out-Of-Memory (OOM)** errors on the UCL CS cluster.
    \begin{alertblock}{The $O(T^2)$ Trap}
        Pre-calculating windows for 50 assets over 30 years created a dataset that exceeded system RAM and VRAM limits during training initialization.
    \end{alertblock}
    \textbf{The Solution:} We refactored to a **Memory-Efficient Lazy Loader**. Windows are sliced on-the-fly from flat 1D arrays using zero-copy stride-tricks, reducing the memory footprint by \textbf{8x}.
\end{frame}

\begin{frame}{Validation: The Breadth-Stability Link}
    \begin{itemize}
        \item \textbf{24-Asset Model:} Achieved 1.27 Sharpe in-sample, but collapsed to -1.00 Out-of-Sample. (The "Narrow Data" Overfit).
        \item \textbf{46-Asset Model:} By doubling the universe, we provided enough cross-sectional variance for the Transformer to learn \textit{regimes} rather than \textit{noise}.
    \end{itemize}
    \textit{Conclusion:} High-capacity models require proportional data breadth. Diversification isn't just for PnL; it's for model regularization.
\end{frame}
